\vfill\pagebreak
\section{Counting: \texttt{for} and ranges}
\label{sec:control:for}

The third loop is for counting. \token{for i in 1 ... 5 \{ ... \}} runs its block
once for each integer from 1 to 5, in order, and inside the block the variable
\token{i} holds the current one. \token{i} is called the \textit{iterator} of
the loop, and \token{1 ... 5} is a \textit{range}.

\begin{lstlisting}[style=coloredverbatim, caption=Counting from 1 to 5, label=lst:control:for]
use std::io;

fn main() {
  for i in 1 ... 5 {
    println(i, " squared is ", i * i);
  }
}
\end{lstlisting}
\vspace{-10pt}%
\begin{lstlisting}[style=bashVerb]
1 squared is 1
2 squared is 4
3 squared is 9
4 squared is 16
5 squared is 25
\end{lstlisting}

\subsection{Ranges}

A range is written with two bounds and one of two operators between them.

\begin{itemize}
\item \token{a ... b}, with three dots, goes from \token{a} to \token{b}, both
  included: \token{1 ... 5} is 1, 2, 3, 4, 5.
\item \token{a .. b}, with two dots, stops just before \token{b}: \token{0 .. 5}
  is 0, 1, 2, 3, 4. It still counts five numbers, and a count from zero is so
  common that this is the form most loops use.
\end{itemize}

The bounds can be any integer expressions, including values read at run time.
They are computed once, before the first iteration. A range whose end comes
before its start, such as \token{10 .. 0}, is empty, and the loop runs zero
times; the next listing relies on that when the user types \token{0}.

\begin{lstlisting}[style=coloredverbatim, caption=A range whose end is read at run time, label=lst:control:for_sum]
use std::io;

fn main() {
  let n = read!i32(ask-> "Up to: ");
  let mut sum = 0;
  for i in 1 ... n {
    sum += i;
  }
  println("1 + 2 + ... + ", n, " = ", sum);
}
\end{lstlisting}
\vspace{-10pt}%
\begin{lstlisting}[style=bashVerb, escapechar=@+]
@+\textcolor{teal!80}{alice@dev:\mtilde}@+$ ./main
Up to: 100
1 + 2 + ... + 100 = 5050
\end{lstlisting}

To count down, call \token{.reverse()} on the range. \token{(1 ... 10).reverse()}
counts from 10 down to 1. The bound that \token{..} leaves out is still the
upper one, so \token{(0 .. 10).reverse()} counts from 9 down to 0, whereas
\token{(0 ... 10).reverse()} starts at 10. To count by steps larger than one,
call \token{.stepBy(k)}, which keeps one number out of every \token{k}. In both
cases the parentheses around the range are needed, so that the call applies to
the whole range rather than to its last bound.

\begin{lstlisting}[style=coloredverbatim, caption=Counting down and by steps, label=lst:control:for_reverse]
use std::io;

fn main() {
  for i in (1 ... 10).reverse() {
    print(i, " ");
  }
  println("liftoff!");

  for i in (0 ... 20).stepBy(5) {
    print(i, " ");
  }
  println();
}
\end{lstlisting}
\vspace{-10pt}%
\begin{lstlisting}[style=bashVerb]
10 9 8 7 6 5 4 3 2 1 liftoff!
0 5 10 15 20
\end{lstlisting}

\token{print}, used here for the first time, is \token{println} without the line
break at the end, so that the numbers stay on one line. A \token{println} with
no argument ends the line.

\subsection{The iterator}

The iterator is a constant. Each iteration gives it the next number of the
range, but the block cannot change it: \token{i = i + 1;} inside the loop is
rejected, because ``i is a value iterator, and cannot be used as a lvalue''. To
skip numbers, use \token{.stepBy} or \token{continue}.

The iterator is also a variable like any other, so the compiler rejects it when
the block never reads it. When only the number of repetitions matters, the
wildcard \token{\_} takes the place of its name. It declares nothing, and is
therefore never unused.

\begin{lstlisting}[style=coloredverbatim, caption=Repeating without reading the iterator, label=lst:control:for_wildcard]
use std::io;

fn main() {
  for _ in 0 .. 3 {
    println("Hip hip hooray!");
  }
}
\end{lstlisting}

\subsection{Loops inside loops}

The block of a loop can contain another loop. The inner loop then runs in full
during each iteration of the outer one. In the next listing the outer loop picks
a row and the inner loop prints the columns of that row.

A \token{break} leaves only the loop it is written in, the innermost one. On
line~7 it ends the current row as soon as the column goes past the row number,
and the outer loop carries on with the next row. That is what gives the output
its triangle shape. \token{continue} likewise affects only the innermost loop.

\begin{lstlisting}[style=coloredverbatim, caption=A \token{break} in a nested loop, label=lst:control:nested]
use std::io;

fn main() {
  for row in 1 ... 4 {
    for col in 1 ... 4 {
      if col > row {
        break;
      }
      print(row * col, " ");
    }
    println();
  }
}
\end{lstlisting}
\vspace{-10pt}%
\begin{lstlisting}[style=bashVerb]
1
2 4
3 6 9
4 8 12 16
\end{lstlisting}

\subsection{Which loop to use}

The three loops overlap, and choosing well makes a program easier to read.

\begin{center}
  \begin{adjustbox}{max width=\linewidth}
    \begin{tabular}{|l|l|}
      \hline
      When & Use\\[0pt]
      \hline
      \hline
      The number of iterations is known before the loop starts & \token{for}\\[0pt]
      The loop goes on as long as a condition, tested first, holds & \token{while}\\[0pt]
      The decision to stop comes after some of the work & \token{loop} and \token{break}\\[0pt]
      \hline
    \end{tabular}
  \end{adjustbox}
\end{center}

A later chapter shows that \token{for} does much more than count: it goes
through the elements of a collection in the same way it goes through the numbers
of a range.
